% !TEX root = ../main.tex
% File: part1/chapters_efficiency/sec2_llm_distillation.tex

\section{Chưng cất tri thức cho LLM}
\label{sec:llm_distillation}

Chưng cất kinh điển (mục~\ref{ssec:knowledge_distillation}) khớp phân phối xác suất của thầy và trò trên \emph{một} nhãn. Với LLM tự hồi quy, ``nhãn'' là cả một \emph{chuỗi} token sinh ra tuần tự, và mỗi token đã sinh lại là điều kiện cho mọi token sau nó. Khác biệt ấy không chỉ làm bài toán lớn hơn về khối lượng tính toán; nó sinh ra ba câu hỏi thiết kế mà chưng cất cho bài toán phân loại không hề có:

\begin{enumerate}
	\item \textbf{Trò học trên những tiền tố nào?} Tiền tố lấy từ văn bản người viết, từ văn bản thầy sinh ra, hay từ chính văn bản trò đang sinh ra trong lúc huấn luyện.
	\item \textbf{Khớp hai phân phối theo độ đo nào?} Tại một tiền tố cho trước, thầy và trò đều cho một phân phối trên toàn từ vựng; so sánh chúng bằng KL xuôi, KL ngược, hay một thỏa hiệp giữa hai chiều.
	\item \textbf{Chưng cất cái gì?} Chỉ đáp án cuối cùng, hay cả chuỗi suy luận dẫn tới đáp án.
\end{enumerate}

Ba câu hỏi này không độc lập với nhau, và chúng cũng không chỉ thuộc về chưng cất: chúng là cùng một bộ câu hỏi mà mọi phương pháp hậu huấn luyện (post-training) đều phải trả lời. Vì vậy, trước khi đi vào từng phương pháp, ta cần một khung phân loại đủ chặt để thấy các phương pháp ấy nằm ở đâu so với nhau.

\subsection{Khung phân loại: nguồn dữ liệu và độ dày tín hiệu}
\label{ssec:post_training_quadrants}

Mọi thuật toán hậu huấn luyện đều nhận vào hai thứ: một tập chuỗi để học, và một tín hiệu nói chuỗi đó tốt hay xấu ở chỗ nào. Hai thứ ấy tạo thành hai trục phân loại độc lập.

\paragraph{Trục thứ nhất: nguồn dữ liệu.} Chuỗi mà ta tính mất mát trên đó được lấy từ phân phối nào?
\begin{itemize}
	\item \textbf{Off-policy.} Chuỗi đến từ một kho cố định, có sẵn trước khi vòng huấn luyện bắt đầu: văn bản người viết, dữ liệu chỉ dẫn, hoặc đầu ra do một mô hình thầy sinh ra từ trước. Tập dữ liệu không đổi khi tham số của trò đổi.
	\item \textbf{On-policy.} Chuỗi do \emph{chính mô hình đang được huấn luyện} sinh ra, tại thời điểm cập nhật hiện tại. Mỗi bước gradient đổi tham số, nên phân phối dữ liệu của bước sau cũng khác bước trước.
\end{itemize}

\paragraph{Trục thứ hai: độ dày tín hiệu.} Trên một chuỗi dài $T$ token, ta nhận được bao nhiêu con số để học?
\begin{itemize}
	\item \textbf{Tín hiệu thưa (sparse).} Một con số cho cả chuỗi: phần thưởng đúng/sai của bộ kiểm chứng, điểm của mô hình thưởng, hoặc một nhãn ``câu A tốt hơn câu B''.
	\item \textbf{Tín hiệu dày (dense).} Một con số cho \emph{mỗi} token: cross-entropy tại từng vị trí, hoặc phân kỳ giữa hai phân phối tại từng vị trí.
\end{itemize}

Cắt hai trục này cho bốn ô, và bốn ô ấy chính là bốn họ phương pháp mà quyển sách đã và đang bàn tới (bảng~\ref{tab:post_training_quadrants}).

\begin{table}[H]
\centering\small
\caption{Bốn ô của hậu huấn luyện, cắt theo nguồn dữ liệu và độ dày tín hiệu.}
\label{tab:post_training_quadrants}
\begin{tabularx}{\textwidth}{@{}>{\raggedright\arraybackslash}p{2.3cm}XX@{}}
\toprule
 & \textbf{Tín hiệu thưa} & \textbf{Tín hiệu dày} \\
\midrule
\textbf{Off-policy} & Tối ưu sở thích ngoại tuyến: DPO và họ (mục~\ref{ssec:dpo}). Một cặp (được chọn, bị loại) đã ghi sẵn, một tín hiệu cho cả câu trả lời. & SFT và chưng cất off-policy: cross-entropy hoặc phân kỳ KL theo từng token trên văn bản cố định. \\
\addlinespace
\textbf{On-policy} & Học tăng cường: PPO (mục~\ref{ssec:rlhf}), GRPO (mục~\ref{ssec:grpo}). Trò tự sinh rollout, nhận một phần thưởng cho cả rollout. & \textbf{Chưng cất on-policy} (mục~\ref{sec:on_policy_distillation}): trò tự sinh rollout, thầy chấm từng token. \\
\bottomrule
\end{tabularx}
\end{table}

\paragraph{Hai trục giải quyết hai vấn đề khác nhau, nên không thay thế được cho nhau.} Trục nguồn dữ liệu quyết định ta có đang tối ưu \emph{đúng đại lượng} hay không. Một hàm mất mát nhỏ trên phân phối $A$ không bảo đảm điều gì về hành vi trên phân phối $B$; mà lúc chạy thật, mô hình luôn đứng trên phân phối của chính nó, không phải trên phân phối của dữ liệu huấn luyện. Trục độ dày quyết định ta học \emph{nhanh} đến đâu: nó đếm số bit thông tin thu được trên mỗi lượt sinh. Một phần thưởng nhị phân cho một chuỗi suy luận dài mười nghìn token là một bit cho mười nghìn quyết định; muốn biết quyết định nào gây ra cái sai, thuật toán phải suy ra từ rất nhiều mẫu khác nhau.

Vì thế ba ô đầu đều hy sinh một thứ. SFT và chưng cất off-policy có tín hiệu dày nhưng đo trên phân phối sai. RL đo trên đúng phân phối nhưng tín hiệu quá thưa. DPO ngoại tuyến hy sinh cả hai, đổi lại là rẻ và ổn định vì không phải sinh văn bản trong vòng huấn luyện.

\paragraph{Ô thứ tư là ô duy nhất không phải chọn một trong hai.} Nếu ta để trò tự sinh chuỗi, rồi hỏi một mô hình thầy ``ở mỗi vị trí, anh sẽ đặt xác suất thế nào'', ta vừa đo trên phân phối trò thật sự sẽ gặp khi chạy, vừa có một con số cho từng token. Đó là chưng cất on-policy.

Cái giá của ô thứ tư cũng rõ ràng, và cần nói ngay để tránh hiểu nó là phương pháp tốt nhất trong mọi trường hợp. Nó đòi một mô hình thầy \emph{đã giỏi sẵn} ở tác vụ đích và chấm được từng token -- điều RL không cần. Chất lượng trò vì thế bị chặn bởi chất lượng thầy, trong khi RL với một bộ kiểm chứng đúng đắn về nguyên tắc có thể vượt mọi mô hình hiện có. Chưng cất on-policy là cách \emph{chuyển} năng lực đã tồn tại ở đâu đó sang một mô hình nhỏ hơn với chi phí thấp; nó không phải cách \emph{tạo ra} năng lực mới.

\begin{keyidea}{Đọc bảng bốn ô như thế nào}
Hai trục trả lời hai câu hỏi khác nhau: \emph{đo trên phân phối nào} (đúng hay sai mục tiêu) và \emph{mỗi lượt sinh cho bao nhiêu bit} (nhanh hay chậm). Phần còn lại của mục này đi theo cột phải, hàng trên: các phương pháp off-policy, tín hiệu dày. Mục~\ref{sec:on_policy_distillation} đi vào ô dưới bên phải.
\end{keyidea}

\subsection{Từ nhãn mềm đến chuỗi: KD gốc và chưng cất mức chuỗi}
\label{ssec:seq_kd}

\paragraph{Nhãn mềm mang gì mà nhãn cứng không có.} Với một nhãn cứng dạng one-hot, cross-entropy chỉ nhìn vào xác suất mà mô hình gán cho lớp đúng. Mọi cách phân bổ phần xác suất còn lại trên các lớp sai đều cho cùng một giá trị mất mát: hàm mất mát không phân biệt được một mô hình cho rằng ảnh con số 2 hơi giống số 3, với một mô hình cho rằng nó hơi giống con bò. Nhưng thầy thì phân biệt được, và nó ghi điều đó vào các xác suất rất nhỏ ở đuôi phân phối. Hinton và cộng sự \cite{hinton2015distilling} phát biểu đúng ý này: ``The relative probabilities of incorrect answers tell us a lot about how the cumbersome model tends to generalize''.

Nói cách khác, nhãn mềm không chỉ nói \emph{đáp án là gì}; nó còn nói \emph{cấu trúc tương tự giữa các đáp án theo cách nhìn của thầy}. Đó là một ràng buộc hình học lên không gian biểu diễn mà nhãn cứng hoàn toàn không có. Bài báo gốc gắn thẳng điều này với hiệu quả dữ liệu: ``When the soft targets have high entropy, they provide much more information per training case than hard targets'', và phương sai của gradient giữa các mẫu cũng nhỏ hơn. Đây chính là bản gốc của lập luận ``tín hiệu dày'' ở trục thứ hai trong bảng~\ref{tab:post_training_quadrants}: cùng một số mẫu, nhãn mềm chở nhiều thông tin hơn hẳn.

\paragraph{Vai trò của nhiệt độ $T$.} Vấn đề là với một thầy đã hội tụ tốt, chính những xác suất mang thông tin ấy lại nhỏ cỡ $10^{-6}$ và gần như không đóng góp gì vào mất mát. Nhiệt độ sinh ra để kéo chúng lên mức nhìn thấy được:
\[
p^{(T)}_i = \frac{\exp(z_i/T)}{\sum_j \exp(z_j/T)} .
\]
Với $T=1$ ta có softmax thường. Tăng $T$ làm phân phối phẳng ra, nâng tỉ trọng tương đối của các lớp đuôi, nên gradient bắt đầu ``nghe thấy'' thông tin ở đuôi. Thầy và trò phải dùng \emph{cùng một} $T$ khi tính mất mát mềm; lúc suy luận thì trò quay về $T=1$.

\paragraph{Vì sao phải nhân lại $T^2$.} Đạo hàm của mất mát mềm theo logit $z_i$ của trò có dạng $\tfrac{1}{T}\big(q^{(T)}_i - p^{(T)}_i\big)$, và bản thân hiệu trong ngoặc cũng co lại theo $1/T$ khi $T$ lớn, vì cả hai phân phối cùng phẳng dần về phân phối đều. Hệ quả là độ lớn gradient của thành phần mềm tỉ lệ với $1/T^2$. Bài báo gốc nói thẳng: ``Since the magnitudes of the gradients produced by the soft targets scale as $1/T^2$'' nên ``it is important to multiply them by $T^2$ when using both hard and soft targets''. Mất mát kết hợp vì vậy có dạng
\[
\mathcal{L} = (1-\alpha)\,\mathcal{L}_{\mathrm{CE}}\big(y,\, q^{(1)}\big)
 \;+\; \alpha\,T^2\,\mathcal{L}_{\mathrm{soft}}\big(p^{(T)},\, q^{(T)}\big).
\]
Lý do thực dụng của hệ số $T^2$: không có nó, mỗi lần đổi $T$ là tỉ lệ đóng góp giữa hai thành phần lại đổi theo, và siêu tham số $\alpha$ phải dò lại từ đầu. Nhân $T^2$ giữ cho hai thành phần cân bằng xấp xỉ như nhau khi $T$ thay đổi, nên $T$ và $\alpha$ có thể dò gần như độc lập.

\paragraph{Word-level KD: áp thẳng ý tưởng trên vào từng vị trí.} Với mô hình chuỗi, cách chuyển trực tiếp nhất là coi mỗi vị trí trong câu là một bài phân loại riêng, trên tiền tố \emph{lấy từ dữ liệu huấn luyện}:
\[
\mathcal{L}_{\text{word}} = \sum_{t} \mathrm{KL}\Big(p_T(\cdot\mid y_{<t}, x)\,\Big\|\, q_S(\cdot\mid y_{<t}, x)\Big).
\]
Mất mát này dày tối đa: mỗi token cho một phân phối đầy đủ trên từ vựng. Nhưng nó đòi hai điều kiện ngặt -- thầy và trò phải dùng chung bộ tách token, và ta phải có quyền đọc logit của thầy.

\paragraph{Sequence-level KD \cite{kim2016sequence}.} Kim và Rush chỉ ra rằng đại lượng ta thật sự quan tâm không phải phân phối tại từng vị trí mà phân phối trên \emph{cả chuỗi}:
\[
\mathcal{L}_{\text{seq}} = -\sum_{\mathbf{t}\in\mathcal{T}} q(\mathbf{t}\mid \mathbf{s})\,\log p(\mathbf{t}\mid \mathbf{s}),
\]
trong đó $\mathcal{T}$ là tập \emph{mọi} chuỗi có thể. Bài báo nhấn mạnh rằng đây là một đại lượng khác hẳn mất mát mức từ, vì ``the sum is over an exponential number of terms''. Tổng đó không tính được, nên các tác giả xấp xỉ bằng cách thay phân phối của thầy bởi \emph{mode} của nó -- và mode thì tìm gần đúng được bằng beam search:
\[
\mathcal{L}_{\text{seq}} \approx -\log p(\mathbf{t} = \hat{\mathbf{y}} \mid \mathbf{s}),
\qquad \hat{\mathbf{y}} = \text{beam-search}(q,\mathbf{s}).
\]
Biện minh cho xấp xỉ này là một giả thiết thực nghiệm về dịch máy: hiệu quả đã biết của beam search khiến các tác giả tin rằng phần lớn khối lượng xác suất của $q$ dồn vào một chuỗi đầu ra duy nhất. Bài báo còn đề xuất một biến thể trung gian, \emph{nội suy mức chuỗi}, trộn mất mát trên câu đích thật với mất mát trên chuỗi trong beam gần câu đích nhất theo BLEU mức câu. Kết quả: mô hình trò chạy nhanh gấp mười lần thầy với mức giảm chất lượng nhỏ, hơn mô hình huấn luyện không chưng cất 4,2 BLEU khi giải mã tham lam và 1,7 BLEU khi dùng beam search; cắt tỉa trọng số thêm cho mô hình ít hơn thầy 13 lần tham số với mức giảm 0,4 BLEU.

\paragraph{Vì sao Seq-KD là nền của ``SFT trên đầu ra của thầy''.} Hãy đọc lại công thức xấp xỉ ở trên. Sau khi thay phân phối của thầy bằng mode của nó, chưng cất mức chuỗi \emph{chính là} cross-entropy thông thường trên một tập văn bản do thầy sinh ra. Không cần logit của thầy, không cần chạy thầy trong vòng huấn luyện, không cần chung bộ tách token, không cần chung kiến trúc. Chỉ cần văn bản. Đó là lý do mọi quy trình dạng ``sinh dữ liệu bằng mô hình lớn rồi SFT mô hình nhỏ trên dữ liệu đó'' -- Alpaca, hay 800 nghìn mẫu chưng cất từ DeepSeek-R1 ở mục~\ref{ssec:r1_distill} -- đều là Seq-KD dưới một cái tên khác, và là lý do nhánh này phổ biến hơn hẳn mọi nhánh còn lại.

Cái mất cũng nằm ngay trong xấp xỉ ấy: ta vứt bỏ toàn bộ phần phân phối nằm ngoài mode. Thầy có thể biết năm cách diễn đạt đều đúng; trò chỉ nhìn thấy một. Với dịch máy, giả thiết ``phần lớn khối lượng dồn vào một chuỗi'' còn hợp lý. Với hội thoại mở, với văn bản sáng tạo, hay với suy luận nhiều bước mà nhiều lối đi khác nhau cùng dẫn tới đáp án đúng, giả thiết đó yếu đi rõ rệt.

\subsection{Lệch phân phối: vì sao huấn luyện ngoài chính sách thì sai dần}
\label{ssec:exposure_bias}

Cả word-level KD lẫn Seq-KD đều huấn luyện trò trên tiền tố \emph{không do chính nó sinh ra}. Đây không phải một chi tiết kỹ thuật nhỏ; nó là lý do tồn tại của toàn bộ nhánh on-policy, nên cần trình bày cho rõ.

\paragraph{Cơ chế.} Trong huấn luyện, ở bước $t$ mô hình được cho sẵn tiền tố đúng $y_{<t}$ và chỉ phải đoán $y_t$. Khi chạy thật, không ai cho nó tiền tố đúng: nó phải tiếp tục từ $\hat{y}_{<t}$ do chính nó vừa sinh ra. Giả sử ở bước nào đó nó chọn nhầm một token. Tiền tố ở bước sau bây giờ là một chuỗi mà \emph{không hề xuất hiện trong dữ liệu huấn luyện}. Trên tiền tố đó, mô hình không có bảo đảm nào về chất lượng dự đoán, nên khả năng sai lại cao hơn mức trung bình, và sai lầm mới lại đẩy tiền tố đi xa hơn nữa khỏi vùng đã học. Vòng lặp này là \textbf{lệch phân phối} (còn gọi là exposure bias, hay compounding error trong học bắt chước).

\paragraph{Gốc lý thuyết: DAgger.} Ross, Gordon và Bagnell \cite{ross2011dagger} đã chứng minh điều này bằng ngôn ngữ của học bắt chước, với một cận rất dễ nhớ. Nếu một chính sách học theo kiểu giám sát thông thường có tỉ lệ sai $\epsilon$ mỗi bước \emph{đo trên phân phối trạng thái của thầy}, thì chi phí của nó trên một bài toán dài $T$ bước bị chặn bởi
\[
J(\pi) \;\le\; J(\pi^\star) + T^2\epsilon ,
\]
tức là sai số tích lũy tăng theo \emph{bình phương} độ dài, chứ không phải tuyến tính. Trực giác đằng sau số mũ hai rất đơn giản: trong $T$ bước có cỡ $T\epsilon$ lần sai, và mỗi lần sai có thể làm hỏng phần còn lại dài cỡ $T$ của quỹ đạo.

Cách sửa mà các tác giả đề xuất, DAgger, là một vòng lặp gộp dữ liệu: cho chính sách \emph{hiện tại} chạy để thu đúng những trạng thái mà nó thật sự ghé qua, hỏi thầy nên làm gì tại chính các trạng thái ấy, gộp các cặp (trạng thái của trò, hành động của thầy) vào tập dữ liệu tích lũy, huấn luyện lại, rồi lặp. Ở vòng đầu, chính sách còn quá kém nên thuật toán cho thầy cầm lái một phần với hệ số trộn $\beta_i$ giảm dần qua các vòng. Chính sách thu được được bảo đảm hoạt động tốt \emph{trên chính phân phối trạng thái mà nó sinh ra}, và cận sai số trở thành tuyến tính theo $T$.

\paragraph{Ánh xạ sang LLM.} Phép tương ứng gần như một-một: trạng thái là tiền tố, hành động là token tiếp theo, chính sách là mô hình trò, thầy là mô hình lớn, và ``hỏi thầy'' là chạy một lượt forward của thầy trên tiền tố do trò sinh ra. Có một khác biệt, và nó làm mọi chuyện dễ hơn nhiều so với bối cảnh robot mà DAgger ra đời: thầy LLM không trả về một hành động duy nhất mà trả về \emph{cả một phân phối trên toàn từ vựng}. Mỗi lần truy vấn thầy vì thế cho nhiều thông tin hơn hẳn một nhãn. Toàn bộ nhánh trình bày ở mục~\ref{sec:on_policy_distillation} có thể đọc như DAgger với một ông thầy giàu thông tin hơn.

\begin{pitfall}{Lệch phân phối không phải là vấn đề chất lượng dữ liệu}
Dữ liệu của thầy có thể hoàn hảo, không nhiễu, không mâu thuẫn, mà trò vẫn sai dần khi chạy. Nguyên nhân không nằm ở \emph{nội dung} dữ liệu mà ở \emph{phân phối} sinh ra nó. Thu thập thêm dữ liệu cùng loại làm giảm $\epsilon$ nhưng không đụng được vào thừa số $T$ thứ hai trong cận $T^2\epsilon$. Muốn bỏ thừa số đó thì phải đổi phân phối lấy mẫu, chứ không phải đổi lượng dữ liệu.
\end{pitfall}

\subsection{Chọn chiều KL trong ngữ cảnh chưng cất}
\label{ssec:kl_direction_distillation}

Tính bất đối xứng của phân kỳ KL và hai hành vi mà nó sinh ra -- bao phủ mode với KL xuôi, bám mode với KL ngược -- đã được trình bày ở mục~\ref{ssec:ly_thuyet_thong_tin}. Ở đây chỉ cần nối nó vào bối cảnh chưng cất, nơi $p$ là thầy và $q_\theta$ là trò.

\paragraph{Khoảng cách năng lực làm chiều KL trở thành vấn đề thật.} Nếu trò đủ lớn để biểu diễn chính xác phân phối của thầy, chiều KL không quan trọng mấy: cực tiểu của cả hai chiều đều là $q_\theta = p$. Chưng cất lại là bài toán trong đó trò \emph{cố tình} nhỏ hơn thầy nhiều lần, nên không có cách nào khớp hoàn hảo, và ta buộc phải quyết định sai ở đâu thì chấp nhận được. KL xuôi phạt vô cùng ở mọi nơi thầy có khối lượng mà trò gán xác suất gần không, nên nó ép trò trải mỏng xác suất ra khắp nơi -- kể cả vào các vùng ``ở giữa các mode'' mà bản thân thầy cũng cho là vô lý. Khi lấy mẫu, trò sẽ có lúc rút trúng chính những vùng đó. KL ngược cho phép trò bỏ hẳn một số mode của thầy mà không bị phạt, đổi lại là mất đa dạng.

\paragraph{Lựa chọn phụ thuộc tác vụ.} Với tác vụ có đáp án đúng (toán, sinh mã, trích xuất thông tin), đa dạng không phải thứ ta cần giữ, và bám mode là lựa chọn đúng. Với sinh dữ liệu tổng hợp, viết sáng tạo, hoặc bất cứ ứng dụng nào mà sự đơn điệu của đầu ra là một lỗi, mất đa dạng là cái giá thật và phải cân nhắc.

\paragraph{Một hệ quả tính toán nối thẳng vào trục thứ nhất.} KL xuôi lấy kỳ vọng theo $p$, tức chỉ cần mẫu của thầy, tức tính được trên một tập dữ liệu tĩnh. KL ngược lấy kỳ vọng theo $q_\theta$, mà $q_\theta$ đổi sau \emph{mỗi} bước cập nhật, nên buộc phải sinh văn bản ngay trong vòng huấn luyện. Nói cách khác, chọn KL ngược \emph{ở mức chuỗi} là đã chọn luôn on-policy ở trục nguồn dữ liệu; hai lựa chọn không độc lập.

Nhưng cần phân biệt cho kỹ, vì đây là chỗ hay lẫn: ta hoàn toàn có thể tính KL ngược \emph{theo từng token}, trên một tiền tố cố định lấy từ dữ liệu. Khi đó chiều KL chỉ nói về cách so sánh hai phân phối \emph{tại một tiền tố cho trước}, còn tiền tố ấy đến từ đâu lại là chuyện của trục nguồn dữ liệu. Chính sự tách bạch này là điều GKD khai thác ở mục~\ref{ssec:gkd}, và nó cho phép dùng KL ngược mà không phải trả giá của policy gradient.

\subsection{MiniLLM: KL ngược và ba kỹ thuật để nó chạy được}
\label{ssec:minillm}

\textbf{MiniLLM} \cite{gu2024minillm} là cách đọc thẳng nhất của lập luận trên: nếu KL ngược là độ đo đúng khi trò nhỏ hơn thầy, thì cứ tối ưu nó ở mức chuỗi,
\[
\theta^\star = \arg\min_\theta \mathrm{KL}\big(q_\theta \,\|\, p_T\big)
= \arg\min_\theta \; \mathbb{E}_{y\sim q_\theta}\left[\log\frac{q_\theta(y\mid x)}{p_T(y\mid x)}\right].
\]
Bài báo mô tả hệ quả mong muốn đúng theo ngôn ngữ của mục trước: trò ``seek the major modes'' của thầy và gán xác suất thấp cho các vùng trống của thầy, thay vì bị ép phủ mọi mẫu của thầy kể cả các mẫu rất khó xảy ra.

\paragraph{Vì sao nó biến thành một bài học tăng cường.} Kỳ vọng lấy theo $y\sim q_\theta$, tức theo chính tham số đang tối ưu, nên không thể lấy đạo hàm trực tiếp qua khâu lấy mẫu. Cách chuẩn là policy gradient. Đặt điểm chất lượng từng bước và tổng tích lũy của nó
\[
r_t = \log\frac{p_T(y_t\mid y_{<t},x)}{q_\theta(y_t\mid y_{<t},x)},
\qquad
R_t = \sum_{t'\ge t} r_{t'},
\]
thì gradient có dạng quen thuộc $-\mathbb{E}_{y\sim q_\theta}\big[\sum_t R_t\,\nabla_\theta \log q_\theta(y_t\mid y_{<t},x)\big]$. Nói cách khác, thầy đóng vai mô hình thưởng, và ta đang chạy RL. Mọi khó khăn của RL vì thế theo vào, và ba kỹ thuật của MiniLLM là ba miếng vá cho đúng ba khó khăn đó.

\begin{enumerate}
	\item \textbf{Phân rã theo từng bước, để giảm phương sai.} Bài báo quan sát rằng ``the error in the front tokens accumulates along the whole sentence'', nên ước lượng Monte Carlo của $R_t$ có phương sai lớn. Tách $R_t = r_t + R_{t+1}$: số hạng $r_t$ chỉ phụ thuộc token tại bước $t$, nên kỳ vọng của nó theo $y_t$ tính được \emph{chính xác} bằng cách cộng trên toàn từ vựng thay vì lấy mẫu. Chỉ phần đuôi $R_{t+1}$ mới cần Monte Carlo. Phương sai giảm, hội tụ nhanh hơn.
	\item \textbf{Lấy mẫu trộn với thầy, để chống khai thác hàm thưởng.} Vì thầy là hàm thưởng, trò có thể tìm ra những chuỗi nhận điểm cao mà thực chất vô nghĩa; bài báo nêu ví dụ các chuỗi lặp cụm từ. Cách sửa là lấy mẫu từ hỗn hợp giữa phân phối của thầy và của trò, với hệ số $\alpha$ điều khiển mức trộn thầy vào ($\alpha = 0{,}2$ trong thí nghiệm). Rollout nhờ vậy ở lại trong vùng mà thầy cho là hợp lý.
	\item \textbf{Chuẩn hóa theo độ dài, để chống thiên lệch câu ngắn.} $R_{t+1}$ là tổng của các log-tỉ số thường mang dấu âm, nên chuỗi càng dài thì tổng càng nhỏ: mục tiêu vô tình thưởng cho câu trả lời ngắn. Cách sửa là lấy trung bình thay cho tổng, chia cho số token còn lại.
\end{enumerate}

\paragraph{Kết quả và cách đọc chúng.} Trên các họ mô hình với quy mô từ 120M đến 13B tham số, MiniLLM cho văn bản ít lệch phân phối hơn, hiệu chỉnh xác suất (calibration) tốt hơn -- đo bằng ECE trên SST2 và BoolQ với LLaMA-7B -- và chất lượng sinh văn bản dài cao hơn các phương pháp chưng cất thông thường, với mức mất đa dạng mà các tác giả mô tả là không đáng kể. Điều đáng rút ra không phải bảng số mà là cấu trúc của phương pháp: ba kỹ thuật ổn định ở trên đều không liên quan gì tới chưng cất, chúng chỉ là giá phải trả cho việc đã chuyển bài toán thành RL. Câu hỏi tự nhiên tiếp theo là liệu có cách nào dùng KL ngược mà không phải trả giá đó không.

\subsection{GKD: tách bạch nguồn chuỗi và độ đo}
\label{ssec:gkd}

\textbf{Generalized Knowledge Distillation (GKD)} \cite{agarwal2024gkd} trả lời đúng câu hỏi vừa nêu, bằng cách tách hai lựa chọn mà MiniLLM gộp làm một thành hai núm vặn độc lập.

\paragraph{Núm thứ nhất: tỉ lệ dữ liệu của trò, $\lambda$.} Ở mỗi bước, với xác suất $\lambda$ thuật toán lấy một lô đầu vào và để \emph{trò} sinh đầu ra; ngược lại nó lấy một lô từ tập dữ liệu cố định (câu đích thật, hoặc đầu ra do thầy sinh sẵn). Tham số $\lambda\in[0,1]$ chính là ``student data fraction'' mà bài báo định nghĩa, tức tọa độ của phương pháp trên trục nguồn dữ liệu ở bảng~\ref{tab:post_training_quadrants}.

\paragraph{Núm thứ hai: phân kỳ, qua Jensen--Shannon tổng quát.} Sau khi đã có lô chuỗi, GKD cực tiểu phân kỳ giữa phân phối token của thầy và của trò tại từng vị trí, với phân kỳ chọn được trong một họ liên tục:
\[
\begin{aligned}
\mathcal{D}_{\mathrm{JSD}(\beta)}(P\|Q) = {}& \beta\,\mathrm{KL}\big(P \,\big\|\, \beta P + (1-\beta)Q\big)\\
&+ (1-\beta)\,\mathrm{KL}\big(Q \,\big\|\, \beta P + (1-\beta)Q\big).
\end{aligned}
\]
Họ này nội suy giữa hai chiều KL với hệ số bị chặn $0<\beta<1$: bài báo nêu rằng gradient của $\mathrm{JSD}(\beta)$ hành xử giống KL xuôi khi $\beta$ gần $0$ và giống KL ngược khi $\beta$ gần $1$. Điểm quan trọng về mặt số học là $\mathrm{JSD}(\beta)$ \emph{bị chặn}: vì cả hai số hạng đều so với phân phối trộn, không có tình huống một token thầy thích mà trò gán xác suất gần không sinh ra hình phạt vô cùng. Đó là lý do nó ổn định hơn KL xuôi thuần túy đúng ở tình huống mà chưng cất luôn rơi vào, tức khi trò không đủ năng lực khớp thầy.

\paragraph{Vì sao GKD không phải là RL.} Đây là điểm kỹ thuật dễ bỏ qua nhất và cũng là điểm phân biệt GKD với MiniLLM. GKD \emph{không} truyền gradient qua khâu lấy mẫu: chuỗi do trò sinh ra, một khi đã sinh, được coi như dữ liệu cố định. Mục tiêu vì thế vẫn là một bài học có giám sát thông thường theo từng token, chỉ khác là tập dữ liệu được làm mới liên tục. Không có hàm thưởng, không có ước lượng gradient phương sai cao, không cần ba miếng vá của MiniLLM. Đổi lại, GKD không tối ưu đúng KL ngược \emph{mức chuỗi} mà tối ưu một phân kỳ mức token trên các chuỗi do trò sinh.

\paragraph{Những gì trở thành trường hợp riêng.} Bài báo phát biểu rằng chưng cất on-policy và chưng cất có giám sát là GKD với KL xuôi và $\lambda$ bằng $1$ và $0$; ImitKD là GKD với KL xuôi và một lịch giảm dần cho $\lambda$; $f$-distill cũng là một trường hợp riêng. Cần nói rõ một điều mà cách trình bày phổ biến hay nói quá: \emph{bài báo GKD không tuyên bố MiniLLM là trường hợp riêng của nó}. Hai phương pháp tối ưu hai đại lượng khác nhau (mức chuỗi so với mức token) bằng hai cơ chế khác nhau (policy gradient so với học có giám sát trên dữ liệu làm mới); GKD được giới thiệu như một phương án đơn giản hơn chứ không phải một phương án bao trùm.

\paragraph{Kết quả và hệ quả kiến trúc.} GKD được đánh giá trên tóm tắt, dịch máy, suy luận số học, và chưng cất không phụ thuộc tác vụ cho tinh chỉnh theo chỉ dẫn. Một hệ quả thực dụng đáng chú ý: vì khâu sinh chuỗi on-policy giống hệt khâu rollout trong RL, hai mục tiêu ghép được vào cùng một vòng huấn luyện, nên chưng cất và tinh chỉnh bằng học tăng cường (ví dụ RLHF, mục~\ref{ssec:rlhf}) không còn là hai giai đoạn tách rời.

\paragraph{Từ bài báo sang quy trình sản xuất.} Dòng mô hình Qwen3 \cite{yang2025qwen3} dùng đúng cấu trúc này để tạo các mô hình nhỏ, với tên gọi \emph{strong-to-weak distillation} và hai giai đoạn nối tiếp. Giai đoạn off-policy chưng cất trên đầu ra của thầy, kết hợp đầu ra sinh ở cả hai chế độ \texttt{/think} và \texttt{/no\_think}, mà báo cáo mô tả là để ``đặt nền vững chắc cho giai đoạn on-policy tiếp theo''. Giai đoạn on-policy cho mô hình trò tự sinh chuỗi rồi căn logit của nó với logit của thầy (Qwen3-32B hoặc Qwen3-235B-A22B) sao cho cực tiểu phân kỳ KL; báo cáo chỉ ghi ``KL divergence'' mà không nói rõ chiều. Chính các con số chi phí của giai đoạn thứ hai là nội dung chính của mục~\ref{sec:on_policy_distillation}.

\subsection{Chưng cất quá trình suy luận}
\label{ssec:rationale_distillation}

Hai mục trên bàn về \emph{cách} khớp phân phối. Mục này bàn về câu hỏi thiết kế thứ ba: chưng cất \emph{cái gì}. Nếu thầy giải được bài toán nhờ một chuỗi suy luận dài, thì học thuộc đáp án cuối là bỏ phí gần hết thông tin mà thầy đã tạo ra.

\paragraph{Distilling Step-by-Step \cite{hsieh2023distilling}.} Thay vì chỉ học đáp án, trò học cả \emph{lý giải} (rationale) do LLM sinh ra bằng Chain-of-Thought, dưới dạng \textbf{học đa nhiệm}:
\[
\mathcal{L} = \mathcal{L}_{\text{label}} + \lambda\, \mathcal{L}_{\text{rationale}},
\]
trong đó cùng một đầu vào được gắn tiền tố \texttt{[label]} để dự đoán đáp án và \texttt{[rationale]} để sinh lý giải. Khi suy luận chỉ cần chạy nhiệm vụ \texttt{[label]}, nên không tốn thêm chi phí sinh lý giải. Lý giải ở đây đóng vai một dạng \emph{giám sát bổ sung giàu thông tin}: nó buộc mô hình nhỏ phải mã hóa được các bước trung gian chứ không chỉ ghi nhớ ánh xạ từ câu hỏi sang đáp án, nên cùng một lượng dữ liệu cho nhiều ràng buộc hơn. Mô hình T5 770M vượt PaLM 540B dùng few-shot prompting trên ANLI chỉ với 80\% dữ liệu huấn luyện, trong khi nhỏ hơn 700 lần.

\paragraph{KARD: chưng cất suy luận có tăng cường tri thức \cite{kang2023kard}.} Với các tác vụ \emph{cần nhiều kiến thức} (y khoa, câu hỏi đa bước), mô hình nhỏ không đủ dung lượng tham số để ghi nhớ mọi kiến thức mà lý giải của thầy dựa vào. Học thuộc lý giải khi đó dạy trò bắt chước hình thức lập luận mà không có dữ kiện để lập luận. KARD kết hợp chưng cất lý giải với truy xuất:
\begin{itemize}
	\item Tinh chỉnh trò sinh lý giải của LLM, \emph{kèm theo} các tài liệu được truy xuất từ kho tri thức bên ngoài.
	\item Huấn luyện một \textbf{bộ xếp hạng lại (reranker)} nơ-ron để tìm các tài liệu thực sự liên quan tới việc sinh lý giải, chứ không chỉ liên quan tới câu hỏi.
\end{itemize}
Kết quả: mô hình T5 250M tham số dùng KARD vượt các mô hình 3B (lớn hơn 12 lần) được tinh chỉnh thông thường trên MedQA-USMLE và StrategyQA. Bài học đáng nhớ hơn con số: \emph{tách kiến thức (đưa vào truy xuất) khỏi kỹ năng suy luận (chưng cất vào tham số)}, tức tinh thần của RAG (chương~\ref{chap:rag}) áp dụng cho mô hình nhỏ.

\subsection{Chưng cất kết hợp học tăng cường: RLAD}
\label{ssec:rlad}

Ở mục~\ref{ssec:r1_distill}, chúng ta thấy chưng cất từ mô hình suy luận lớn hiệu quả hơn RL trực tiếp trên mô hình nhỏ. Đọc theo bảng~\ref{tab:post_training_quadrants}, câu hỏi tiếp theo là làm sao dùng cùng lúc hai nguồn tín hiệu ở hai ô khác nhau: phần thưởng thưa nhưng đúng đắn của bộ kiểm chứng, và tín hiệu dày nhưng bị chặn trên của thầy. Cách đơn giản nhất là thêm một số hạng KL giữa trò và thầy vào mục tiêu RL. Zhang và cộng sự \cite{zhang2026rlad} chỉ ra hai vấn đề của cách này: dữ liệu của thầy bị \textbf{lệch phân phối} so với các chuỗi do trò sinh ra trong RL, và số hạng KL \textbf{xung đột mục tiêu} với việc cực đại phần thưởng -- kéo trò về phía thầy ngay cả ở những chỗ mà làm vậy khiến phần thưởng giảm.

\textbf{RLAD (RL-Aware Distillation)} đề xuất \emph{bắt chước có chọn lọc}: chỉ kéo trò về phía thầy khi điều đó thực sự cải thiện bước cập nhật chính sách hiện tại. Thành phần chính, \textbf{Trust Region Ratio Distillation (TRRD)}, thay số hạng KL thầy--trò bằng một mục tiêu dạng tỷ số xác suất kiểu PPO/GRPO, được neo vào một \emph{hỗn hợp giữa thầy và chính sách cũ}. Nhờ vậy việc chưng cất vừa nhận biết advantage (chỉ học từ thầy ở nơi có lợi) vừa bị giới hạn trong vùng tin cậy (không nhảy quá xa khỏi chính sách hiện tại). Trên các benchmark suy luận logic và toán học, RLAD vượt cả chưng cất ngoại tuyến, GRPO tiêu chuẩn, và chưng cất on-policy dựa trên KL.

\begin{table}[H]
\centering\small
\caption{Bản đồ các phương pháp chưng cất cho LLM, đọc theo hai trục của bảng~\ref{tab:post_training_quadrants}.}
\label{tab:llm_distillation_map}
\begin{tabularx}{\textwidth}{@{}>{\raggedright\arraybackslash}p{2.9cm}XX@{}}
\toprule
\textbf{Phương pháp} & \textbf{Trò học từ chuỗi nào, tín hiệu gì} & \textbf{Vì sao thiết kế như vậy} \\
\midrule
Seq-KD / SFT trên dữ liệu thầy & Chuỗi của thầy (mode), cross-entropy & Thay phân phối chuỗi của thầy bằng mode nên chỉ cần văn bản; đơn giản nhất, nhưng lệch phân phối và mất mọi thứ ngoài mode \\
\addlinespace
Word-level KD & Tiền tố cố định, KL xuôi từng token & Tín hiệu dày nhất có thể, nhưng đòi chung bộ tách token và quyền đọc logit thầy \\
\addlinespace
MiniLLM & Chuỗi của trò, KL ngược mức chuỗi & Bám mode khi trò nhỏ hơn thầy; phải dùng policy gradient nên cần ba kỹ thuật ổn định \\
\addlinespace
GKD & Trộn chuỗi cố định và của trò theo $\lambda$, $\mathrm{JSD}(\beta)$ từng token & Tách nguồn chuỗi khỏi độ đo; phân kỳ bị chặn nên ổn định; không truyền gradient qua khâu lấy mẫu nên không phải RL \\
\addlinespace
Distilling Step-by-Step & Đáp án và lý giải của thầy & Lý giải là giám sát phụ, ràng buộc các bước trung gian chứ không chỉ kết quả \\
\addlinespace
KARD & Lý giải của thầy kèm tài liệu truy xuất & Mô hình nhỏ không đủ chỗ chứa kiến thức, nên đẩy kiến thức ra ngoài tham số \\
\addlinespace
RLAD & Chuỗi của trò, phần thưởng và thầy có chọn lọc & Gộp tín hiệu thưa đúng đắn với tín hiệu dày bị chặn trên, trong vùng tin cậy \\
\bottomrule
\end{tabularx}
\end{table}

\begin{assess}{Phần còn thiếu của bức tranh}
Bảng~\ref{tab:llm_distillation_map} còn trống đúng một ô: trò tự sinh chuỗi, thầy chấm từng token, và mục tiêu là một phân kỳ từng token chứ không phải một phần thưởng cho cả chuỗi. GKD đã mở đường tới đó nhưng trình bày nó như một lựa chọn trong nhiều lựa chọn; những gì học được từ đó sau này -- cách đặt nó vào đúng hạ tầng RL sẵn có, chi phí thực tế so với RL, và các điều kiện để nó không thất bại -- là nội dung của mục~\ref{sec:on_policy_distillation}.
\end{assess}
